Work rate formula:  Rate \times Time = Work

For identical workers:  R_{total} = n \cdot R_{one}

Man-hours:  \frac{M_1 T_1}{W_1} = \frac{M_2 T_2}{W_2}

Example (2026 \#10): 5 mules eat 3 bales in 6 hours
Rate per mule = \frac{3}{5 \cdot 6} = \frac{1}{10}  bales/hour
Time for 5 mules to eat 4 bales = \frac{4}{5 \cdot 1/10} = 8  hours

Mason example (2018 \#16): 6 masons tile 540 ft^2  in 4 hours
Rate = \frac{540}{6 \cdot 4} = 22.5  ft^2/mason/hour
9 masons, 1620 ft^2: time = \frac{1620}{9 \cdot 22.5} = 8  hours

Painting (2020 \#16): 2 men paint 3 houses in 4 weeks
Rate = \frac{3}{2 \cdot 4} = \frac{3}{8}  houses/man/week
1 house, 2 men:  t = \frac{1}{2 \cdot 3/8} = \frac{4}{3}  weeks

Mixture problem (State 2022 \#14): 7000 L at 20\% OJ
Drain  x  liters, replace with 80\% OJ:
0.20(7000-x) + 0.80x = 0.25(7000)
1400 - 0.20x + 0.80x = 1750
0.60x = 350 \rightarrow x = \frac{1750}{3}

Mixing concentrations (2026 \#4): 1 L of 10\% + 4 L of  k\% = 5 L of 8\%
0.10(1) + \frac{k}{100}(4) = 0.08(5)
0.10 + 0.04k = 0.40
0.04k = 0.30 \rightarrow k = 7.5
